%% Generated by Sphinx.
\def\sphinxdocclass{report}
\documentclass[letterpaper,10pt,polish]{sphinxmanual}
\ifdefined\pdfpxdimen
   \let\sphinxpxdimen\pdfpxdimen\else\newdimen\sphinxpxdimen
\fi \sphinxpxdimen=.75bp\relax
\ifdefined\pdfimageresolution
    \pdfimageresolution= \numexpr \dimexpr1in\relax/\sphinxpxdimen\relax
\fi
%% let collapsible pdf bookmarks panel have high depth per default
\PassOptionsToPackage{bookmarksdepth=5}{hyperref}

\PassOptionsToPackage{booktabs}{sphinx}
\PassOptionsToPackage{colorrows}{sphinx}

\PassOptionsToPackage{warn}{textcomp}
\usepackage[utf8]{inputenc}
\ifdefined\DeclareUnicodeCharacter
% support both utf8 and utf8x syntaxes
  \ifdefined\DeclareUnicodeCharacterAsOptional
    \def\sphinxDUC#1{\DeclareUnicodeCharacter{"#1}}
  \else
    \let\sphinxDUC\DeclareUnicodeCharacter
  \fi
  \sphinxDUC{00A0}{\nobreakspace}
  \sphinxDUC{2500}{\sphinxunichar{2500}}
  \sphinxDUC{2502}{\sphinxunichar{2502}}
  \sphinxDUC{2514}{\sphinxunichar{2514}}
  \sphinxDUC{251C}{\sphinxunichar{251C}}
  \sphinxDUC{2572}{\textbackslash}
\fi
\usepackage{cmap}
\usepackage[T1]{fontenc}
\usepackage{amsmath,amssymb,amstext}
\usepackage{babel}



\usepackage{tgtermes}
\usepackage{tgheros}
\renewcommand{\ttdefault}{txtt}



\usepackage[Sonny]{fncychap}
\ChNameVar{\Large\normalfont\sffamily}
\ChTitleVar{\Large\normalfont\sffamily}
\usepackage{sphinx}

\fvset{fontsize=auto}
\usepackage{geometry}


% Include hyperref last.
\usepackage{hyperref}
% Fix anchor placement for figures with captions.
\usepackage{hypcap}% it must be loaded after hyperref.
% Set up styles of URL: it should be placed after hyperref.
\urlstyle{same}

\addto\captionspolish{\renewcommand{\contentsname}{Contents:}}

\usepackage{sphinxmessages}
\setcounter{tocdepth}{1}



\title{Kopie zapasowe i odzyskiwanie danych}
\date{18 cze 2026}
\release{0.0.1}
\author{Gal 1}
\newcommand{\sphinxlogo}{\vbox{}}
\renewcommand{\releasename}{Wydanie}
\makeindex
\begin{document}

\ifdefined\shorthandoff
  \ifnum\catcode`\=\string=\active\shorthandoff{=}\fi
  \ifnum\catcode`\"=\active\shorthandoff{"}\fi
\fi

\pagestyle{empty}
\sphinxmaketitle
\pagestyle{plain}
\sphinxtableofcontents
\pagestyle{normal}
\phantomsection\label{\detokenize{index::doc}}


\sphinxAtStartPar
Add your content using \sphinxcode{\sphinxupquote{reStructuredText}} syntax. See the
\sphinxhref{https://www.sphinx-doc.org/en/master/usage/restructuredtext/index.html}{reStructuredText}
documentation for details.

\sphinxstepscope

\sphinxstepscope


\chapter{Dokumentacja bazy danych \textendash{} Sklep elektroniczny}
\label{\detokenize{rozdzial_3/index:dokumentacja-bazy-danych-sklep-elektroniczny}}\label{\detokenize{rozdzial_3/index::doc}}
\sphinxAtStartPar
\sphinxstylestrong{Autorzy:} Norbert Antonovitch, Michał Bednarczyk, Jan Balazs de Borbatviz


\section{Wybór zagadnienia i identyfikacja danych}
\label{\detokenize{rozdzial_3/index:wybor-zagadnienia-i-identyfikacja-danych}}
\sphinxAtStartPar
Tematem projektu jest podstawowa baza danych obsługująca sklep internetowy ze sprzętem elektronicznym.
Baza danych zawiera następujące grupy danych:
\begin{itemize}
\item {} 
\sphinxAtStartPar
\sphinxstylestrong{Klienci:} unikalny identyfikator, imię, nazwisko, adres e\sphinxhyphen{}mail, numer telefonu oraz podstawowe dane adresowe do wysyłki: ulica, miasto, kod pocztowy.

\item {} 
\sphinxAtStartPar
\sphinxstylestrong{Produkty:} unikalny identyfikator, nazwa sprzętu, cena jednostkowa oraz przypisana kategoria.

\item {} 
\sphinxAtStartPar
\sphinxstylestrong{Kategorie:} identyfikator kategorii oraz jej nazwa, na przykład Laptopy lub Smartfony.

\item {} 
\sphinxAtStartPar
\sphinxstylestrong{Zamówienia:} identyfikator zamówienia, data złożenia, status, na przykład nowe lub zrealizowane, oraz powiązanie z konkretnym klientem.

\item {} 
\sphinxAtStartPar
\sphinxstylestrong{Szczegóły zamówienia:} pozycje przypisane do konkretnego zamówienia, określające jaki produkt został kupiony, w jakiej ilości oraz po jakiej cenie.

\end{itemize}


\section{Opis procesów i towarzyszących im danych}
\label{\detokenize{rozdzial_3/index:opis-procesow-i-towarzyszacych-im-danych}}\begin{enumerate}
\sphinxsetlistlabels{\arabic}{enumi}{enumii}{}{.}%
\item {} 
\sphinxAtStartPar
\sphinxstylestrong{Zarządzanie produktami:} sklep oferuje asortyment przypisany do konkretnych kategorii, na przykład laptopy i smartfony. Każdy produkt ma przypisaną nazwę, cenę oraz identyfikator producenta.

\item {} 
\sphinxAtStartPar
\sphinxstylestrong{Zarządzanie użytkownikami:} klienci rejestrują się w systemie, podając podstawowe dane, takie jak imię, nazwisko i e\sphinxhyphen{}mail, oraz pojedynczy główny adres zamieszkania wykorzystywany do wysyłki.

\item {} 
\sphinxAtStartPar
\sphinxstylestrong{Realizacja zamówień:} zalogowany klient składa zamówienie na wybrane produkty. System tworzy nagłówek zamówienia zawierający datę i status oraz przypisuje do niego poszczególne produkty wraz z ich ilością.

\end{enumerate}


\section{Prototypowy plik JSON}
\label{\detokenize{rozdzial_3/index:prototypowy-plik-json}}
\sphinxAtStartPar
Poniżej przedstawiono prototypowy plik w formacie JSON, zawierający jeden pełny dokument reprezentujący złożone zamówienie. Pozwala to zweryfikować kompletność gromadzonych i przetwarzanych informacji przed przejściem do modelowania konceptualnego.

\begin{sphinxVerbatim}[commandchars=\\\{\}]
\PYG{p}{\PYGZob{}}
\PYG{+w}{  }\PYG{n+nt}{\PYGZdq{}id\PYGZus{}zamowienia\PYGZdq{}}\PYG{p}{:}\PYG{+w}{ }\PYG{l+m+mi}{1}\PYG{p}{,}
\PYG{+w}{  }\PYG{n+nt}{\PYGZdq{}data\PYGZus{}zlozenia\PYGZdq{}}\PYG{p}{:}\PYG{+w}{ }\PYG{l+s+s2}{\PYGZdq{}2026\PYGZhy{}05\PYGZhy{}14T10:15:30Z\PYGZdq{}}\PYG{p}{,}
\PYG{+w}{  }\PYG{n+nt}{\PYGZdq{}status\PYGZdq{}}\PYG{p}{:}\PYG{+w}{ }\PYG{l+s+s2}{\PYGZdq{}Nowe\PYGZdq{}}\PYG{p}{,}
\PYG{+w}{  }\PYG{n+nt}{\PYGZdq{}klient\PYGZdq{}}\PYG{p}{:}\PYG{+w}{ }\PYG{p}{\PYGZob{}}
\PYG{+w}{    }\PYG{n+nt}{\PYGZdq{}id\PYGZus{}klienta\PYGZdq{}}\PYG{p}{:}\PYG{+w}{ }\PYG{l+m+mi}{4096}\PYG{p}{,}
\PYG{+w}{    }\PYG{n+nt}{\PYGZdq{}imie\PYGZdq{}}\PYG{p}{:}\PYG{+w}{ }\PYG{l+s+s2}{\PYGZdq{}Jan\PYGZdq{}}\PYG{p}{,}
\PYG{+w}{    }\PYG{n+nt}{\PYGZdq{}nazwisko\PYGZdq{}}\PYG{p}{:}\PYG{+w}{ }\PYG{l+s+s2}{\PYGZdq{}Kowalski\PYGZdq{}}\PYG{p}{,}
\PYG{+w}{    }\PYG{n+nt}{\PYGZdq{}email\PYGZdq{}}\PYG{p}{:}\PYG{+w}{ }\PYG{l+s+s2}{\PYGZdq{}student01@example.com\PYGZdq{}}\PYG{p}{,}
\PYG{+w}{    }\PYG{n+nt}{\PYGZdq{}telefon\PYGZdq{}}\PYG{p}{:}\PYG{+w}{ }\PYG{l+s+s2}{\PYGZdq{}123456789\PYGZdq{}}\PYG{p}{,}
\PYG{+w}{    }\PYG{n+nt}{\PYGZdq{}ulica\PYGZdq{}}\PYG{p}{:}\PYG{+w}{ }\PYG{l+s+s2}{\PYGZdq{}Armii Krajowej 78\PYGZdq{}}\PYG{p}{,}
\PYG{+w}{    }\PYG{n+nt}{\PYGZdq{}kod\PYGZus{}pocztowy\PYGZdq{}}\PYG{p}{:}\PYG{+w}{ }\PYG{l+s+s2}{\PYGZdq{}58\PYGZhy{}300\PYGZdq{}}\PYG{p}{,}
\PYG{+w}{    }\PYG{n+nt}{\PYGZdq{}miasto\PYGZdq{}}\PYG{p}{:}\PYG{+w}{ }\PYG{l+s+s2}{\PYGZdq{}Walbrzych\PYGZdq{}}
\PYG{+w}{  }\PYG{p}{\PYGZcb{},}
\PYG{+w}{  }\PYG{n+nt}{\PYGZdq{}pozycje\PYGZdq{}}\PYG{p}{:}\PYG{+w}{ }\PYG{p}{[}
\PYG{+w}{    }\PYG{p}{\PYGZob{}}
\PYG{+w}{      }\PYG{n+nt}{\PYGZdq{}id\PYGZus{}produktu\PYGZdq{}}\PYG{p}{:}\PYG{+w}{ }\PYG{l+m+mi}{10543}\PYG{p}{,}
\PYG{+w}{      }\PYG{n+nt}{\PYGZdq{}nazwa\PYGZdq{}}\PYG{p}{:}\PYG{+w}{ }\PYG{l+s+s2}{\PYGZdq{}Laptop ASUS\PYGZdq{}}\PYG{p}{,}
\PYG{+w}{      }\PYG{n+nt}{\PYGZdq{}kategoria\PYGZdq{}}\PYG{p}{:}\PYG{+w}{ }\PYG{l+s+s2}{\PYGZdq{}Laptopy\PYGZdq{}}\PYG{p}{,}
\PYG{+w}{      }\PYG{n+nt}{\PYGZdq{}ilosc\PYGZdq{}}\PYG{p}{:}\PYG{+w}{ }\PYG{l+m+mi}{1}\PYG{p}{,}
\PYG{+w}{      }\PYG{n+nt}{\PYGZdq{}cena\PYGZus{}historyczna\PYGZdq{}}\PYG{p}{:}\PYG{+w}{ }\PYG{l+m+mf}{6299.99}
\PYG{+w}{    }\PYG{p}{\PYGZcb{}}
\PYG{+w}{  }\PYG{p}{]}
\PYG{p}{\PYGZcb{}}
\end{sphinxVerbatim}


\subsection{Dobór typów danych}
\label{\detokenize{rozdzial_3/index:dobor-typow-danych}}

\begin{savenotes}\sphinxattablestart
\sphinxthistablewithglobalstyle
\centering
\sphinxcapstartof{table}
\sphinxthecaptionisattop
\sphinxcaption{Mapowanie typów danych dla SQLite i PostgreSQL}\label{\detokenize{rozdzial_3/index:id1}}
\sphinxaftertopcaption
\begin{tabulary}{\linewidth}[t]{TTT}
\sphinxtoprule
\sphinxstyletheadfamily 
\sphinxAtStartPar
Atrybut danych
&\sphinxstyletheadfamily 
\sphinxAtStartPar
SQLite
&\sphinxstyletheadfamily 
\sphinxAtStartPar
PostgreSQL
\\
\sphinxmidrule
\sphinxtableatstartofbodyhook
\sphinxAtStartPar
Identyfikatory (\sphinxcode{\sphinxupquote{id\_...}})
&
\sphinxAtStartPar
\sphinxcode{\sphinxupquote{INTEGER}}
&
\sphinxAtStartPar
\sphinxcode{\sphinxupquote{SERIAL}} lub \sphinxcode{\sphinxupquote{INTEGER}}
\\
\sphinxhline
\sphinxAtStartPar
Napisy krótkie (\sphinxcode{\sphinxupquote{imie}}, \sphinxcode{\sphinxupquote{nazwisko}}, \sphinxcode{\sphinxupquote{miasto}})
&
\sphinxAtStartPar
\sphinxcode{\sphinxupquote{TEXT}}
&
\sphinxAtStartPar
\sphinxcode{\sphinxupquote{VARCHAR(50)}}
\\
\sphinxhline
\sphinxAtStartPar
Napisy długie (\sphinxcode{\sphinxupquote{ulica}}, \sphinxcode{\sphinxupquote{email}}, \sphinxcode{\sphinxupquote{nazwa}})
&
\sphinxAtStartPar
\sphinxcode{\sphinxupquote{TEXT}}
&
\sphinxAtStartPar
\sphinxcode{\sphinxupquote{VARCHAR(255)}}
\\
\sphinxhline
\sphinxAtStartPar
Stałe ciągi znaków (\sphinxcode{\sphinxupquote{telefon}}, \sphinxcode{\sphinxupquote{kod\_pocztowy}})
&
\sphinxAtStartPar
\sphinxcode{\sphinxupquote{TEXT}}
&
\sphinxAtStartPar
\sphinxcode{\sphinxupquote{VARCHAR(15)}}
\\
\sphinxhline
\sphinxAtStartPar
Cena (\sphinxcode{\sphinxupquote{cena\_bazowa}}, \sphinxcode{\sphinxupquote{cena\_historyczna}})
&
\sphinxAtStartPar
\sphinxcode{\sphinxupquote{FLOAT}}
&
\sphinxAtStartPar
\sphinxcode{\sphinxupquote{NUMERIC(10,2)}}
\\
\sphinxhline
\sphinxAtStartPar
Ilość (\sphinxcode{\sphinxupquote{ilosc}})
&
\sphinxAtStartPar
\sphinxcode{\sphinxupquote{INTEGER}}
&
\sphinxAtStartPar
\sphinxcode{\sphinxupquote{SMALLINT}}
\\
\sphinxhline
\sphinxAtStartPar
Data i czas (\sphinxcode{\sphinxupquote{data\_zlozenia}})
&
\sphinxAtStartPar
\sphinxcode{\sphinxupquote{TEXT}}
&
\sphinxAtStartPar
\sphinxcode{\sphinxupquote{TIMESTAMP}}
\\
\sphinxbottomrule
\end{tabulary}
\sphinxtableafterendhook\par
\sphinxattableend\end{savenotes}


\section{Modelowanie konceptualne i logiczne}
\label{\detokenize{rozdzial_3/index:modelowanie-konceptualne-i-logiczne}}

\subsection{Model koncepcyjny}
\label{\detokenize{rozdzial_3/index:model-koncepcyjny}}
\sphinxAtStartPar
\sphinxstylestrong{1. Encje mocne}
\begin{itemize}
\item {} 
\sphinxAtStartPar
\sphinxstylestrong{Klient:} klucz główny (\sphinxcode{\sphinxupquote{id\_klienta}}), imię, nazwisko, e\sphinxhyphen{}mail, telefon, ulica, miasto, kod pocztowy.

\item {} 
\sphinxAtStartPar
\sphinxstylestrong{Produkt:} klucz główny (\sphinxcode{\sphinxupquote{id\_produktu}}), nazwa, \sphinxcode{\sphinxupquote{cena\_bazowa}}.

\item {} 
\sphinxAtStartPar
\sphinxstylestrong{Kategoria:} klucz główny (\sphinxcode{\sphinxupquote{id\_kategorii}}), nazwa.

\item {} 
\sphinxAtStartPar
\sphinxstylestrong{Zamówienie:} klucz główny (\sphinxcode{\sphinxupquote{id\_zamowienia}}), \sphinxcode{\sphinxupquote{data\_zlozenia}}, status.

\end{itemize}

\sphinxAtStartPar
\sphinxstylestrong{2. Encje słabe}
\begin{itemize}
\item {} 
\sphinxAtStartPar
\sphinxstylestrong{Szczegóły zamówienia:} klucz obcy (\sphinxcode{\sphinxupquote{id\_zamowienia}}), klucz obcy (\sphinxcode{\sphinxupquote{id\_produktu}}), ilość, \sphinxcode{\sphinxupquote{cena\_historyczna}}. Jej identyfikacja opiera się wyłącznie na kluczach obcych pochodzących z encji silnych.

\end{itemize}

\sphinxAtStartPar
\sphinxstylestrong{3. Opis związków}

\sphinxAtStartPar
Zdefiniowano następujące relacje biznesowe między encjami:
\begin{itemize}
\item {} 
\sphinxAtStartPar
\sphinxstylestrong{Klient \textendash{} Zamówienie (1:N):} klient składa zamówienie. Jeden klient może posiadać wiele zamówień, ale jedno konkretne zamówienie jest przypisane dokładnie do jednego klienta.

\item {} 
\sphinxAtStartPar
\sphinxstylestrong{Kategoria \textendash{} Produkt (1:N):} kategoria grupuje produkty. Kategoria może zawierać wiele produktów, ale dany produkt należy do jednej kategorii.

\item {} 
\sphinxAtStartPar
\sphinxstylestrong{Zamówienie \textendash{} Produkt (M:N):} zamówienie zawiera produkty. Pojedyncze zamówienie obejmuje wiele produktów, a dany produkt może pojawić się w wielu zamówieniach.

\end{itemize}

\sphinxAtStartPar
\sphinxstylestrong{4. Związki niepoprawne}

\sphinxAtStartPar
Wyeliminowano bezpośrednią relację między Klientem a Produktem, aby uniknąć pułapki połączeń. Bezpośrednie powiązanie gubi kontekst transakcji, na przykład datę i ilość zakupionego towaru, dlatego poprawna relacja musi zawsze przechodzić przez encję Zamówienie.


\subsection{Model koncepcyjny \textendash{} notacja Chena}
\label{\detokenize{rozdzial_3/index:model-koncepcyjny-notacja-chena}}\begin{quote}

\noindent\sphinxincludegraphics[width=0.800\linewidth]{{schemat_chen}.png}
\end{quote}


\subsection{Model logiczny i normalizacja}
\label{\detokenize{rozdzial_3/index:model-logiczny-i-normalizacja}}
\sphinxAtStartPar
Główne działania podjęte podczas normalizacji:
\begin{enumerate}
\sphinxsetlistlabels{\arabic}{enumi}{enumii}{}{.}%
\item {} 
\sphinxAtStartPar
\sphinxstylestrong{I Postać Normalna (1NF):} dane adresowe klienta, takie jak ulica, miasto i kod pocztowy, zostały wyodrębnione jako pojedyncze pola.

\item {} 
\sphinxAtStartPar
\sphinxstylestrong{II Postać Normalna (2NF):} zlikwidowano relację wiele\sphinxhyphen{}do\sphinxhyphen{}wielu (N:M) pomiędzy Zamówieniem a Produktem, wstawiając encję asocjacyjną \sphinxcode{\sphinxupquote{Szczegóły\_zamówienia}}. Jej klucz główny składa się z identyfikatora zamówienia oraz identyfikatora produktu.

\item {} 
\sphinxAtStartPar
\sphinxstylestrong{III Postać Normalna (3NF):} usunięto zależności przechodnie. Nazwa kategorii została wydzielona do osobnej tabeli Kategorie. W tabeli Produkty przechowywany jest jedynie klucz obcy do kategorii. Ponadto zapewniono zapisywanie historycznej ceny w \sphinxcode{\sphinxupquote{Szczegółach\_zamówienia}}, aby modyfikacja cennika nie fałszowała archiwalnych rachunków.

\end{enumerate}


\subsection{Schemat notacji IE po normalizacji}
\label{\detokenize{rozdzial_3/index:schemat-notacji-ie-po-normalizacji}}\begin{quote}

\noindent\sphinxincludegraphics[width=0.800\linewidth]{{schemat_ie}.png}
\end{quote}

\sphinxstepscope


\chapter{4. Definiowanie bazy danych i wprowadzanie danych do bazy}
\label{\detokenize{rozdzial_4/index:definiowanie-bazy-danych-i-wprowadzanie-danych-do-bazy}}\label{\detokenize{rozdzial_4/index::doc}}
\sphinxAtStartPar
\sphinxstylestrong{Autorzy:} Norbert Antonovitch, Michał Bednarczyk, Jan Balazs de Borbatviz


\section{4.1. Definicja bazy danych w wariantach PostgreSQL i SQLite}
\label{\detokenize{rozdzial_4/index:definicja-bazy-danych-w-wariantach-postgresql-i-sqlite}}
\sphinxAtStartPar
W projekcie zastosowano dwa warianty definicji schematu bazy danych: dla PostgreSQL oraz dla SQLite. Oba warianty odwzorowują tę samą strukturę logiczną, jednak różnią się doborem typów danych i składnią definicji kluczy głównych oraz obcych.

\sphinxAtStartPar
W wariancie PostgreSQL przyjęto rygorystyczne typowanie danych. Identyfikatory zdefiniowano z użyciem typu \sphinxcode{\sphinxupquote{SERIAL}}, pola tekstowe opisano typem \sphinxcode{\sphinxupquote{VARCHAR}} o ograniczonej długości, natomiast wartości finansowe zapisano w typie \sphinxcode{\sphinxupquote{NUMERIC(10,2)}}, zapewniającym odpowiednią precyzję obliczeń.

\begin{sphinxVerbatim}[commandchars=\\\{\}]
\PYG{c+c1}{\PYGZhy{}\PYGZhy{} Wariant PostgreSQL}
\PYG{k}{CREATE}\PYG{+w}{ }\PYG{k}{TABLE}\PYG{+w}{ }\PYG{n}{kategorie}\PYG{+w}{ }\PYG{p}{(}
\PYG{+w}{   }\PYG{n}{id\PYGZus{}kategorii}\PYG{+w}{ }\PYG{n+nb}{SERIAL}\PYG{+w}{ }\PYG{k}{PRIMARY}\PYG{+w}{ }\PYG{k}{KEY}\PYG{p}{,}
\PYG{+w}{   }\PYG{n}{nazwa}\PYG{+w}{ }\PYG{n+nb}{VARCHAR}\PYG{p}{(}\PYG{l+m+mi}{100}\PYG{p}{)}\PYG{+w}{ }\PYG{k}{NOT}\PYG{+w}{ }\PYG{k}{NULL}
\PYG{p}{)}\PYG{p}{;}

\PYG{k}{CREATE}\PYG{+w}{ }\PYG{k}{TABLE}\PYG{+w}{ }\PYG{n}{klienci}\PYG{+w}{ }\PYG{p}{(}
\PYG{+w}{   }\PYG{n}{id\PYGZus{}klienta}\PYG{+w}{ }\PYG{n+nb}{SERIAL}\PYG{+w}{ }\PYG{k}{PRIMARY}\PYG{+w}{ }\PYG{k}{KEY}\PYG{p}{,}
\PYG{+w}{   }\PYG{n}{imie}\PYG{+w}{ }\PYG{n+nb}{VARCHAR}\PYG{p}{(}\PYG{l+m+mi}{50}\PYG{p}{)}\PYG{+w}{ }\PYG{k}{NOT}\PYG{+w}{ }\PYG{k}{NULL}\PYG{p}{,}
\PYG{+w}{   }\PYG{n}{nazwisko}\PYG{+w}{ }\PYG{n+nb}{VARCHAR}\PYG{p}{(}\PYG{l+m+mi}{50}\PYG{p}{)}\PYG{+w}{ }\PYG{k}{NOT}\PYG{+w}{ }\PYG{k}{NULL}\PYG{p}{,}
\PYG{+w}{   }\PYG{n}{email}\PYG{+w}{ }\PYG{n+nb}{VARCHAR}\PYG{p}{(}\PYG{l+m+mi}{255}\PYG{p}{)}\PYG{+w}{ }\PYG{k}{UNIQUE}\PYG{+w}{ }\PYG{k}{NOT}\PYG{+w}{ }\PYG{k}{NULL}\PYG{p}{,}
\PYG{+w}{   }\PYG{n}{telefon}\PYG{+w}{ }\PYG{n+nb}{VARCHAR}\PYG{p}{(}\PYG{l+m+mi}{15}\PYG{p}{)}\PYG{p}{,}
\PYG{+w}{   }\PYG{n}{ulica}\PYG{+w}{ }\PYG{n+nb}{VARCHAR}\PYG{p}{(}\PYG{l+m+mi}{255}\PYG{p}{)}\PYG{p}{,}
\PYG{+w}{   }\PYG{n}{miasto}\PYG{+w}{ }\PYG{n+nb}{VARCHAR}\PYG{p}{(}\PYG{l+m+mi}{50}\PYG{p}{)}\PYG{p}{,}
\PYG{+w}{   }\PYG{n}{kod\PYGZus{}pocztowy}\PYG{+w}{ }\PYG{n+nb}{VARCHAR}\PYG{p}{(}\PYG{l+m+mi}{15}\PYG{p}{)}
\PYG{p}{)}\PYG{p}{;}

\PYG{k}{CREATE}\PYG{+w}{ }\PYG{k}{TABLE}\PYG{+w}{ }\PYG{n}{produkty}\PYG{+w}{ }\PYG{p}{(}
\PYG{+w}{   }\PYG{n}{id\PYGZus{}produktu}\PYG{+w}{ }\PYG{n+nb}{SERIAL}\PYG{+w}{ }\PYG{k}{PRIMARY}\PYG{+w}{ }\PYG{k}{KEY}\PYG{p}{,}
\PYG{+w}{   }\PYG{n}{id\PYGZus{}kategorii}\PYG{+w}{ }\PYG{n+nb}{INTEGER}\PYG{+w}{ }\PYG{k}{NOT}\PYG{+w}{ }\PYG{k}{NULL}\PYG{p}{,}
\PYG{+w}{   }\PYG{n}{nazwa}\PYG{+w}{ }\PYG{n+nb}{VARCHAR}\PYG{p}{(}\PYG{l+m+mi}{255}\PYG{p}{)}\PYG{+w}{ }\PYG{k}{NOT}\PYG{+w}{ }\PYG{k}{NULL}\PYG{p}{,}
\PYG{+w}{   }\PYG{n}{cena\PYGZus{}bazowa}\PYG{+w}{ }\PYG{n+nb}{NUMERIC}\PYG{p}{(}\PYG{l+m+mi}{10}\PYG{p}{,}\PYG{l+m+mi}{2}\PYG{p}{)}\PYG{+w}{ }\PYG{k}{NOT}\PYG{+w}{ }\PYG{k}{NULL}\PYG{p}{,}
\PYG{+w}{   }\PYG{k}{CONSTRAINT}\PYG{+w}{ }\PYG{n}{fk\PYGZus{}kategoria}
\PYG{+w}{      }\PYG{k}{FOREIGN}\PYG{+w}{ }\PYG{k}{KEY}\PYG{+w}{ }\PYG{p}{(}\PYG{n}{id\PYGZus{}kategorii}\PYG{p}{)}\PYG{+w}{ }\PYG{k}{REFERENCES}\PYG{+w}{ }\PYG{n}{kategorie}\PYG{p}{(}\PYG{n}{id\PYGZus{}kategorii}\PYG{p}{)}
\PYG{p}{)}\PYG{p}{;}

\PYG{k}{CREATE}\PYG{+w}{ }\PYG{k}{TABLE}\PYG{+w}{ }\PYG{n}{zamowienia}\PYG{+w}{ }\PYG{p}{(}
\PYG{+w}{   }\PYG{n}{id\PYGZus{}zamowienia}\PYG{+w}{ }\PYG{n+nb}{SERIAL}\PYG{+w}{ }\PYG{k}{PRIMARY}\PYG{+w}{ }\PYG{k}{KEY}\PYG{p}{,}
\PYG{+w}{   }\PYG{n}{id\PYGZus{}klienta}\PYG{+w}{ }\PYG{n+nb}{INTEGER}\PYG{+w}{ }\PYG{k}{NOT}\PYG{+w}{ }\PYG{k}{NULL}\PYG{p}{,}
\PYG{+w}{   }\PYG{n}{data\PYGZus{}zlozenia}\PYG{+w}{ }\PYG{k}{TIMESTAMP}\PYG{+w}{ }\PYG{k}{NOT}\PYG{+w}{ }\PYG{k}{NULL}\PYG{p}{,}
\PYG{+w}{   }\PYG{n}{status}\PYG{+w}{ }\PYG{n+nb}{VARCHAR}\PYG{p}{(}\PYG{l+m+mi}{50}\PYG{p}{)}\PYG{+w}{ }\PYG{k}{NOT}\PYG{+w}{ }\PYG{k}{NULL}\PYG{p}{,}
\PYG{+w}{   }\PYG{k}{CONSTRAINT}\PYG{+w}{ }\PYG{n}{fk\PYGZus{}klient}
\PYG{+w}{      }\PYG{k}{FOREIGN}\PYG{+w}{ }\PYG{k}{KEY}\PYG{+w}{ }\PYG{p}{(}\PYG{n}{id\PYGZus{}klienta}\PYG{p}{)}\PYG{+w}{ }\PYG{k}{REFERENCES}\PYG{+w}{ }\PYG{n}{klienci}\PYG{p}{(}\PYG{n}{id\PYGZus{}klienta}\PYG{p}{)}
\PYG{p}{)}\PYG{p}{;}

\PYG{k}{CREATE}\PYG{+w}{ }\PYG{k}{TABLE}\PYG{+w}{ }\PYG{n}{szczegoly\PYGZus{}zamowienia}\PYG{+w}{ }\PYG{p}{(}
\PYG{+w}{   }\PYG{n}{id\PYGZus{}zamowienia}\PYG{+w}{ }\PYG{n+nb}{INTEGER}\PYG{+w}{ }\PYG{k}{NOT}\PYG{+w}{ }\PYG{k}{NULL}\PYG{p}{,}
\PYG{+w}{   }\PYG{n}{id\PYGZus{}produktu}\PYG{+w}{ }\PYG{n+nb}{INTEGER}\PYG{+w}{ }\PYG{k}{NOT}\PYG{+w}{ }\PYG{k}{NULL}\PYG{p}{,}
\PYG{+w}{   }\PYG{n}{ilosc}\PYG{+w}{ }\PYG{n+nb}{SMALLINT}\PYG{+w}{ }\PYG{k}{NOT}\PYG{+w}{ }\PYG{k}{NULL}\PYG{+w}{ }\PYG{k}{CHECK}\PYG{+w}{ }\PYG{p}{(}\PYG{n}{ilosc}\PYG{+w}{ }\PYG{o}{\PYGZgt{}}\PYG{+w}{ }\PYG{l+m+mi}{0}\PYG{p}{)}\PYG{p}{,}
\PYG{+w}{   }\PYG{n}{cena\PYGZus{}historyczna}\PYG{+w}{ }\PYG{n+nb}{NUMERIC}\PYG{p}{(}\PYG{l+m+mi}{10}\PYG{p}{,}\PYG{l+m+mi}{2}\PYG{p}{)}\PYG{+w}{ }\PYG{k}{NOT}\PYG{+w}{ }\PYG{k}{NULL}\PYG{p}{,}
\PYG{+w}{   }\PYG{k}{PRIMARY}\PYG{+w}{ }\PYG{k}{KEY}\PYG{+w}{ }\PYG{p}{(}\PYG{n}{id\PYGZus{}zamowienia}\PYG{p}{,}\PYG{+w}{ }\PYG{n}{id\PYGZus{}produktu}\PYG{p}{)}\PYG{p}{,}
\PYG{+w}{   }\PYG{k}{CONSTRAINT}\PYG{+w}{ }\PYG{n}{fk\PYGZus{}zamowienie}
\PYG{+w}{      }\PYG{k}{FOREIGN}\PYG{+w}{ }\PYG{k}{KEY}\PYG{+w}{ }\PYG{p}{(}\PYG{n}{id\PYGZus{}zamowienia}\PYG{p}{)}\PYG{+w}{ }\PYG{k}{REFERENCES}\PYG{+w}{ }\PYG{n}{zamowienia}\PYG{p}{(}\PYG{n}{id\PYGZus{}zamowienia}\PYG{p}{)}\PYG{p}{,}
\PYG{+w}{   }\PYG{k}{CONSTRAINT}\PYG{+w}{ }\PYG{n}{fk\PYGZus{}produkt}
\PYG{+w}{      }\PYG{k}{FOREIGN}\PYG{+w}{ }\PYG{k}{KEY}\PYG{+w}{ }\PYG{p}{(}\PYG{n}{id\PYGZus{}produktu}\PYG{p}{)}\PYG{+w}{ }\PYG{k}{REFERENCES}\PYG{+w}{ }\PYG{n}{produkty}\PYG{p}{(}\PYG{n}{id\PYGZus{}produktu}\PYG{p}{)}
\PYG{p}{)}\PYG{p}{;}
\end{sphinxVerbatim}

\sphinxAtStartPar
Wariant SQLite został dostosowany do prostszego modelu typów danych oferowanego przez ten silnik. Typ \sphinxcode{\sphinxupquote{SERIAL}} zastąpiono konstrukcją \sphinxcode{\sphinxupquote{INTEGER PRIMARY KEY AUTOINCREMENT}}, pola tekstowe oraz daty zapisano jako \sphinxcode{\sphinxupquote{TEXT}}, natomiast wartości liczbowe o charakterze ciągłym opisano typem \sphinxcode{\sphinxupquote{FLOAT}}.

\begin{sphinxVerbatim}[commandchars=\\\{\}]
\PYG{c+c1}{\PYGZhy{}\PYGZhy{} Wariant SQLite}
\PYG{k}{CREATE}\PYG{+w}{ }\PYG{k}{TABLE}\PYG{+w}{ }\PYG{n}{kategorie}\PYG{+w}{ }\PYG{p}{(}
\PYG{+w}{   }\PYG{n}{id\PYGZus{}kategorii}\PYG{+w}{ }\PYG{n+nb}{INTEGER}\PYG{+w}{ }\PYG{k}{PRIMARY}\PYG{+w}{ }\PYG{k}{KEY}\PYG{+w}{ }\PYG{n}{AUTOINCREMENT}\PYG{p}{,}
\PYG{+w}{   }\PYG{n}{nazwa}\PYG{+w}{ }\PYG{n+nb}{TEXT}\PYG{+w}{ }\PYG{k}{NOT}\PYG{+w}{ }\PYG{k}{NULL}
\PYG{p}{)}\PYG{p}{;}

\PYG{k}{CREATE}\PYG{+w}{ }\PYG{k}{TABLE}\PYG{+w}{ }\PYG{n}{klienci}\PYG{+w}{ }\PYG{p}{(}
\PYG{+w}{   }\PYG{n}{id\PYGZus{}klienta}\PYG{+w}{ }\PYG{n+nb}{INTEGER}\PYG{+w}{ }\PYG{k}{PRIMARY}\PYG{+w}{ }\PYG{k}{KEY}\PYG{+w}{ }\PYG{n}{AUTOINCREMENT}\PYG{p}{,}
\PYG{+w}{   }\PYG{n}{imie}\PYG{+w}{ }\PYG{n+nb}{TEXT}\PYG{+w}{ }\PYG{k}{NOT}\PYG{+w}{ }\PYG{k}{NULL}\PYG{p}{,}
\PYG{+w}{   }\PYG{n}{nazwisko}\PYG{+w}{ }\PYG{n+nb}{TEXT}\PYG{+w}{ }\PYG{k}{NOT}\PYG{+w}{ }\PYG{k}{NULL}\PYG{p}{,}
\PYG{+w}{   }\PYG{n}{email}\PYG{+w}{ }\PYG{n+nb}{TEXT}\PYG{+w}{ }\PYG{k}{UNIQUE}\PYG{+w}{ }\PYG{k}{NOT}\PYG{+w}{ }\PYG{k}{NULL}\PYG{p}{,}
\PYG{+w}{   }\PYG{n}{telefon}\PYG{+w}{ }\PYG{n+nb}{TEXT}\PYG{p}{,}
\PYG{+w}{   }\PYG{n}{ulica}\PYG{+w}{ }\PYG{n+nb}{TEXT}\PYG{p}{,}
\PYG{+w}{   }\PYG{n}{miasto}\PYG{+w}{ }\PYG{n+nb}{TEXT}\PYG{p}{,}
\PYG{+w}{   }\PYG{n}{kod\PYGZus{}pocztowy}\PYG{+w}{ }\PYG{n+nb}{TEXT}
\PYG{p}{)}\PYG{p}{;}

\PYG{k}{CREATE}\PYG{+w}{ }\PYG{k}{TABLE}\PYG{+w}{ }\PYG{n}{produkty}\PYG{+w}{ }\PYG{p}{(}
\PYG{+w}{   }\PYG{n}{id\PYGZus{}produktu}\PYG{+w}{ }\PYG{n+nb}{INTEGER}\PYG{+w}{ }\PYG{k}{PRIMARY}\PYG{+w}{ }\PYG{k}{KEY}\PYG{+w}{ }\PYG{n}{AUTOINCREMENT}\PYG{p}{,}
\PYG{+w}{   }\PYG{n}{id\PYGZus{}kategorii}\PYG{+w}{ }\PYG{n+nb}{INTEGER}\PYG{+w}{ }\PYG{k}{NOT}\PYG{+w}{ }\PYG{k}{NULL}\PYG{p}{,}
\PYG{+w}{   }\PYG{n}{nazwa}\PYG{+w}{ }\PYG{n+nb}{TEXT}\PYG{+w}{ }\PYG{k}{NOT}\PYG{+w}{ }\PYG{k}{NULL}\PYG{p}{,}
\PYG{+w}{   }\PYG{n}{cena\PYGZus{}bazowa}\PYG{+w}{ }\PYG{n+nb}{FLOAT}\PYG{+w}{ }\PYG{k}{NOT}\PYG{+w}{ }\PYG{k}{NULL}\PYG{p}{,}
\PYG{+w}{   }\PYG{k}{FOREIGN}\PYG{+w}{ }\PYG{k}{KEY}\PYG{+w}{ }\PYG{p}{(}\PYG{n}{id\PYGZus{}kategorii}\PYG{p}{)}\PYG{+w}{ }\PYG{k}{REFERENCES}\PYG{+w}{ }\PYG{n}{kategorie}\PYG{p}{(}\PYG{n}{id\PYGZus{}kategorii}\PYG{p}{)}
\PYG{p}{)}\PYG{p}{;}

\PYG{k}{CREATE}\PYG{+w}{ }\PYG{k}{TABLE}\PYG{+w}{ }\PYG{n}{zamowienia}\PYG{+w}{ }\PYG{p}{(}
\PYG{+w}{   }\PYG{n}{id\PYGZus{}zamowienia}\PYG{+w}{ }\PYG{n+nb}{INTEGER}\PYG{+w}{ }\PYG{k}{PRIMARY}\PYG{+w}{ }\PYG{k}{KEY}\PYG{+w}{ }\PYG{n}{AUTOINCREMENT}\PYG{p}{,}
\PYG{+w}{   }\PYG{n}{id\PYGZus{}klienta}\PYG{+w}{ }\PYG{n+nb}{INTEGER}\PYG{+w}{ }\PYG{k}{NOT}\PYG{+w}{ }\PYG{k}{NULL}\PYG{p}{,}
\PYG{+w}{   }\PYG{n}{data\PYGZus{}zlozenia}\PYG{+w}{ }\PYG{n+nb}{TEXT}\PYG{+w}{ }\PYG{k}{NOT}\PYG{+w}{ }\PYG{k}{NULL}\PYG{p}{,}
\PYG{+w}{   }\PYG{n}{status}\PYG{+w}{ }\PYG{n+nb}{TEXT}\PYG{+w}{ }\PYG{k}{NOT}\PYG{+w}{ }\PYG{k}{NULL}\PYG{p}{,}
\PYG{+w}{   }\PYG{k}{FOREIGN}\PYG{+w}{ }\PYG{k}{KEY}\PYG{+w}{ }\PYG{p}{(}\PYG{n}{id\PYGZus{}klienta}\PYG{p}{)}\PYG{+w}{ }\PYG{k}{REFERENCES}\PYG{+w}{ }\PYG{n}{klienci}\PYG{p}{(}\PYG{n}{id\PYGZus{}klienta}\PYG{p}{)}
\PYG{p}{)}\PYG{p}{;}

\PYG{k}{CREATE}\PYG{+w}{ }\PYG{k}{TABLE}\PYG{+w}{ }\PYG{n}{szczegoly\PYGZus{}zamowienia}\PYG{+w}{ }\PYG{p}{(}
\PYG{+w}{   }\PYG{n}{id\PYGZus{}zamowienia}\PYG{+w}{ }\PYG{n+nb}{INTEGER}\PYG{+w}{ }\PYG{k}{NOT}\PYG{+w}{ }\PYG{k}{NULL}\PYG{p}{,}
\PYG{+w}{   }\PYG{n}{id\PYGZus{}produktu}\PYG{+w}{ }\PYG{n+nb}{INTEGER}\PYG{+w}{ }\PYG{k}{NOT}\PYG{+w}{ }\PYG{k}{NULL}\PYG{p}{,}
\PYG{+w}{   }\PYG{n}{ilosc}\PYG{+w}{ }\PYG{n+nb}{INTEGER}\PYG{+w}{ }\PYG{k}{NOT}\PYG{+w}{ }\PYG{k}{NULL}\PYG{+w}{ }\PYG{k}{CHECK}\PYG{+w}{ }\PYG{p}{(}\PYG{n}{ilosc}\PYG{+w}{ }\PYG{o}{\PYGZgt{}}\PYG{+w}{ }\PYG{l+m+mi}{0}\PYG{p}{)}\PYG{p}{,}
\PYG{+w}{   }\PYG{n}{cena\PYGZus{}historyczna}\PYG{+w}{ }\PYG{n+nb}{FLOAT}\PYG{+w}{ }\PYG{k}{NOT}\PYG{+w}{ }\PYG{k}{NULL}\PYG{p}{,}
\PYG{+w}{   }\PYG{k}{PRIMARY}\PYG{+w}{ }\PYG{k}{KEY}\PYG{+w}{ }\PYG{p}{(}\PYG{n}{id\PYGZus{}zamowienia}\PYG{p}{,}\PYG{+w}{ }\PYG{n}{id\PYGZus{}produktu}\PYG{p}{)}\PYG{p}{,}
\PYG{+w}{   }\PYG{k}{FOREIGN}\PYG{+w}{ }\PYG{k}{KEY}\PYG{+w}{ }\PYG{p}{(}\PYG{n}{id\PYGZus{}zamowienia}\PYG{p}{)}\PYG{+w}{ }\PYG{k}{REFERENCES}\PYG{+w}{ }\PYG{n}{zamowienia}\PYG{p}{(}\PYG{n}{id\PYGZus{}zamowienia}\PYG{p}{)}\PYG{p}{,}
\PYG{+w}{   }\PYG{k}{FOREIGN}\PYG{+w}{ }\PYG{k}{KEY}\PYG{+w}{ }\PYG{p}{(}\PYG{n}{id\PYGZus{}produktu}\PYG{p}{)}\PYG{+w}{ }\PYG{k}{REFERENCES}\PYG{+w}{ }\PYG{n}{produkty}\PYG{p}{(}\PYG{n}{id\PYGZus{}produktu}\PYG{p}{)}
\PYG{p}{)}\PYG{p}{;}
\end{sphinxVerbatim}


\section{4.2. Wybór mechanizmów wsadowego wprowadzania danych}
\label{\detokenize{rozdzial_4/index:wybor-mechanizmow-wsadowego-wprowadzania-danych}}
\sphinxAtStartPar
W celu sprawnego zainicjowania bazy danymi testowymi zrezygnowano z ręcznego przygotowywania pojedynczych instrukcji \sphinxcode{\sphinxupquote{INSERT}} na rzecz automatyzacji realizowanej w języku Python. Dane wejściowe przygotowano w postaci pliku \sphinxcode{\sphinxupquote{JSON}}, co umożliwiło jego łatwe przetwarzanie i przenoszenie pomiędzy wariantami bazy.

\sphinxAtStartPar
Do komunikacji z obiema bazami wykorzystano odpowiednie interfejsy programistyczne: moduł \sphinxcode{\sphinxupquote{sqlite3}} dla SQLite oraz bibliotekę ORM \sphinxcode{\sphinxupquote{SQLAlchemy}} dla PostgreSQL. Istotnym elementem implementacji było zastosowanie wstawiania wsadowego przy użyciu funkcji \sphinxcode{\sphinxupquote{executemany()}}. Mechanizm ten pozwala przesyłać wiele rekordów w ramach jednego wywołania, co ogranicza narzut komunikacyjny i skraca czas zasilania bazy danymi.

\sphinxAtStartPar
Zastosowanie podejścia wsadowego ma szczególne znaczenie w przypadku większej liczby rekordów, ponieważ redukuje liczbę pojedynczych operacji na serwerze i pozwala lepiej wykorzystać mechanizm transakcyjny systemu zarządzania bazą danych. W praktyce przekłada się to na większą wydajność niż przy ręcznym wykonywaniu kolejnych poleceń \sphinxcode{\sphinxupquote{INSERT}}.


\section{4.3. Komentarz do procesu wprowadzania danych}
\label{\detokenize{rozdzial_4/index:komentarz-do-procesu-wprowadzania-danych}}
\sphinxAtStartPar
Proces zasilania bazy danych musiał uwzględniać zależności wynikające z więzów integralności referencyjnej. Z tego względu kolejność importu danych nie mogła być przypadkowa i musiała odzwierciedlać strukturę relacyjną modelu.

\sphinxAtStartPar
Kolejność wprowadzania danych była następująca:
\begin{enumerate}
\sphinxsetlistlabels{\arabic}{enumi}{enumii}{}{.}%
\item {} 
\sphinxAtStartPar
W pierwszej kolejności zaimportowano dane do tabel niezależnych, czyli \sphinxcode{\sphinxupquote{kategorie}} oraz \sphinxcode{\sphinxupquote{klienci}}.

\item {} 
\sphinxAtStartPar
Następnie zasilono tabelę \sphinxcode{\sphinxupquote{produkty}}. Operacja ta była możliwa dopiero po wcześniejszym utworzeniu kategorii, ponieważ każdy produkt musi być przypisany do istniejącej kategorii.

\item {} 
\sphinxAtStartPar
W kolejnym kroku wprowadzono rekordy do tabeli \sphinxcode{\sphinxupquote{zamowienia}}, przypisując je do istniejących klientów.

\item {} 
\sphinxAtStartPar
Na końcu uzupełniono tabelę \sphinxcode{\sphinxupquote{szczegoly\_zamowienia}}, która pełni funkcję encji asocjacyjnej i łączy identyfikatory zamówień oraz produktów. Wymaga to wcześniejszego istnienia obu powiązanych rekordów.

\end{enumerate}

\sphinxAtStartPar
Takie uporządkowanie procesu importu danych wynika bezpośrednio z konstrukcji schematu relacyjnego i stanowi warunek poprawnego zachowania integralności bazy. W praktyce oznacza to, że dane nadrzędne muszą zostać zapisane przed danymi zależnymi.



\renewcommand{\indexname}{Indeks}
\printindex
\end{document}